\documentclass[10pt,a4paper]{amsart}
\usepackage[margin=19mm]{geometry}
\usepackage{amsmath,amssymb,mathtools}
\usepackage[scr=rsfs,cal=euler]{mathalfa}
\usepackage{xfrac}
\usepackage[unicode,hidelinks,bookmarks=false]{hyperref}
\newcommand{\ie}{i.e.\ }
\newcommand{\cf}{cf.\ }
\newcommand{\resp}{respectively\ }
\newcommand{\Subsection}{\subsection}
\providecommand{\nobreakdash}{\nobreak\hbox{-}}
\setlength{\emergencystretch}{2em}
\multlinegap0pt
\usepackage{bm}
\usepackage{bbm}
\usepackage{dsfont}
\usepackage{xspace}
\usepackage{cancel}
\usepackage{mathtools}
\usepackage{amsthm}
\makeatletter
\@ifundefined{theorem}{\newtheorem{theorem}{Theorem}}{}
\@ifundefined{lemma}{\newtheorem{lemma}[theorem]{Lemma}}{}
\makeatother
\title[On ascent sequences avoiding 021 and a pattern of length fou]{On ascent sequences avoiding 021 and a pattern of length four}
\author{Toufik Mansour, Mark Shattuck}
\date{Source: arXiv:2507.17947v1 (CC BY 4.0)}
\begin{document}
\maketitle

\noindent\emph{Source and licence.} On ascent sequences avoiding 021 and a pattern of length four. Version arXiv:2507.17947v1 (CC BY 4.0), distributed under the Creative Commons Attribution 4.0 International licence, as recorded in the arXiv licence metadata for this exact version and shown on the abstract page https://arxiv.org/abs/2507.17947v1. Full rights notice: licenses/arxiv-2507.17947-CC-BY-4.0.txt.\par\smallskip

\noindent\textit{Editorial scope.} Counted unit: the Theorem whose source label is \texttt{1200th} in the article (source0.tex lines 385--390, source label \texttt{1200th}), delivered together with the article's own complete original proof of it (source0.tex lines 391--479, one \texttt{proof} environment of 89 lines). No mathematics is written, repaired or re-derived: statement and proof are the article's own text line for line. Required context retained uncounted: 1000lem (lines 171--173). Every definition, display and label of those lines is retained unchanged. Omitted from this package and not redistributed: 1--170, 174--384, 480--748. Nothing inside a retained excerpt is omitted or altered: every retained body is delivered exactly as the article's lines are written, so no omission receipt of any kind accompanies this package. Citations: the retained excerpts carry no citation command at all, so no bibliography entry of the article is transcribed and no reference is redistributed. Reference adaptation: none was needed and none was made; every reference inside the delivered unit resolves inside it, so no unresolved reference or citation is produced. Printed slips kept literal: no printed word, symbol, hypothesis or number of the counted statement or of its proof was altered. Layout and class: the article's own class is \texttt{article} with packages including amssymb, pstricks, latexsym, amsmath, amsthm; the wrapper is the lane's pinned \texttt{amsart} editorial wrapper at 10pt on A4, which restates the theorem-like environments the retained text needs and adds the article's own macro definitions that the retained text needs (none), with extra wrapper packages (bm, bbm, dsfont, xspace, cancel, mathtools). The delivered numbering is the unit's own. Assets: no retained span contains a figure environment, a graphics-inclusion command, a PDF-page inclusion, a table or a TikZ picture; no image file is copied into this package, the package declares no asset dependency and no image file is redistributed with it. Licence: version arXiv:2507.17947v1 of this article is distributed under the Creative Commons Attribution 4.0 International licence, as recorded in the arXiv licence metadata for this exact version and shown on the abstract page https://arxiv.org/abs/2507.17947v1. A full rights notice is delivered at licenses/arxiv-2507.17947-CC-BY-4.0.txt. Characters: every retained excerpt is pure ASCII, so the delivered engine, XeLaTeX with its default Latin Modern text font, reports no missing character.
\begin{lemma}\label{1000lem}
If $r \geq0$, then $h_r(x)=\frac{x}{1-2x}\left(\frac{x}{1-x}\right)^r$.
\end{lemma}

\begin{theorem}\label{1200th}
If $\tau \in \{1200,1220,1230\}$, then
\begin{equation}\label{1200the1}
f_\tau(x)=\frac{1-7x+17x^2-16x^3+5x^4-x^5}{(1-2x)(1-3x+x^2)^2}.
\end{equation}
\end{theorem}

\begin{proof}
We first find $f_{1230}$. The binary members, together with $\varepsilon$, are accounted for by $\frac{1-x}{1-2x}$, so assume henceforth $\pi \in B(1230)$ contains a letter great than 1.  First suppose $\pi$ can be decomposed as
$$\pi=0^{a_0}1^{b_1}0^{a_1}\cdots 1^{b_\ell}0^{a_\ell}\pi', \qquad \ell \geq 1,$$
where $a_i>0$ for $0 \leq i \leq \ell-1$ with $a _\ell \geq0$, $b_i>0$ for all $i$ and $\pi'$ is a nonempty sequence on $\{0,u\}$ for some $u>1$ that starts with $u$.  Note $\pi$ an ascent sequence implies $2 \leq u \leq \ell+1$. Considering all possible $\ell$, we have that $\pi$ of the stated form are enumerated by
\begin{align*}
\sum_{\ell \geq 1}\frac{x^{2\ell}}{(1-x)^{2\ell+1}}\cdot\frac{\ell x}{1-2x}&=\frac{x^3}{(1-x)^3(1-2x)}\sum_{\ell \geq1}\ell\left(\frac{x^2}{(1-x)^2}\right)^{\ell-1}=\frac{x^3}{(1-x)^3(1-2x)}\cdot\frac{1}{\left(1-\frac{x^2}{(1-x)^2}\right)^2}\\
&=\frac{x^3(1-x)}{(1-2x)^3}.
\end{align*}

So assume $\pi$ contains at least three distinct positive letters.  Then we have
$$\pi=0^{a_0}1^{b_1}0^{a_1}\cdots1^{b_\ell}0^{a_{\ell}}u^{c_1}0^{d_1}\cdots u^{c_p}0^{d_p}\pi', \qquad \ell, p \geq 1,$$
where the $a_i$, $b_i$ and $u$ are as before, $c_i>0$ for all $i$, $d_i>0$ for $1 \leq i \leq p-1$ with $d_p\geq0$ and $\pi' \neq \varepsilon$ starts with $v$ for some $v>u$.  Note that there are $\ell+p$ ascents occurring strictly prior to the first $v$ and hence $v \in [u+1,\ell+p+1]$.  Then $\pi$ avoiding $021$ and $1230$ implies $\pi'$ is nondecreasing and hence is accounted for by $h_{\ell+p+1-v}(x)$, as it contains at most $\ell+p+1-v$ jumps, where $h_i(x)$ is as in Lemma \ref{1000lem}.  Conversely, any $\pi$ of the stated form with $\pi'$ as described is seen to belong to $B(1230)$. Considering all possible $\ell$, $p$, $u$ and $v$, we have that $\pi$ of the form above are enumerated by
$$\sum_{\ell \geq1}\sum_{p \geq1}\sum_{u=2}^{\ell+1}\sum_{v=u+1}^{\ell+p+1}\frac{x^{2\ell}}{(1-x)^{2\ell+1}}\cdot\frac{x^{2p-1}}{(1-x)^{2p}}\cdot h_{\ell+p+1-v}(x).$$

This last sum may be simplified as follows:
\begin{align*}
&\sum_{\ell \geq1}\sum_{p \geq1}\sum_{u=2}^{\ell+1}\sum_{v=u+1}^{\ell+p+1}\frac{x^{2\ell+2p-1}}{(1-x)^{2\ell+2p+1}}\cdot\frac{x}{1-2x}\left(\frac{1-x}{1-2x}\right)^{\ell+p+1-v}\\
&=\frac{1}{(1-2x)^2}\sum_{\ell \geq1}\sum_{p \geq1}\left(\frac{x^2}{(1-x)(1-2x)}\right)^{\ell+p}\sum_{u=2}^{\ell+1}\sum_{v=u+1}^{\ell+p+1}\left(\frac{1-2x}{1-x}\right)^v\\
&=\frac{1}{x(1-2x)}\sum_{\ell \geq1}\sum_{p \geq1}\left(\frac{x^2}{(1-x)(1-2x)}\right)^{\ell+p}\sum_{u=2}^{\ell+1}\left(\left(\frac{1-2x}{1-x}\right)^u-\left(\frac{1-2x}{1-x}\right)^{\ell+p+1}\right)\qquad\qquad\qquad\qquad\quad\\
&=\frac{1}{x(1-2x)}\sum_{\ell \geq1}\sum_{p \geq1}\left(\frac{x^2}{(1-x)(1-2x)}\right)^{\ell+p}\left(\frac{\left(\frac{1-2x}{1-x}\right)^2-\left(\frac{1-2x}{1-x}\right)^{\ell+2}}{1-\frac{1-2x}{1-x}}-\ell\left(\frac{1-2x}{1-x}\right)^{\ell+p+1}\right)\\
&=\frac{1-2x}{x^2(1-x)}\sum_{\ell \geq1}\sum_{p \geq1}\left(\frac{x^2}{(1-x)(1-2x)}\right)^{\ell+p}\left(1-\left(\frac{1-2x}{1-x}\right)^\ell\right)-\frac{1}{x(1-x)}\sum_{\ell\geq1}\sum_{p\geq1}\ell\left(\frac{x^2}{(1-x)^2}\right)^{\ell+p}\\
&=\frac{1-2x}{x^2(1-x)}\cdot\frac{x^4}{(1-3x+x^2)^2}-\frac{1-2x}{x^2(1-x)}\cdot\frac{x^4}{(1-2x)(1-3x+x^2)}-\frac{1}{x(1-x)}\cdot\frac{x^4(1-x)^2}{(1-2x)^3}\\
&=\frac{x^3}{(1-3x+x^2)^2}-\frac{x^3(1-x)}{(1-2x)^3}.
\end{align*}
Combining the last case with the two prior implies
$$f_{1230}=\frac{1-x}{1-2x}+\frac{x^3}{(1-3x+x^2)^2}=\frac{1-7x+17x^2-16x^3+5x^4-x^5}{(1-2x)(1-3x+x^2)^2},$$
as desired.

We now find $f_{1200}$. The binary sequences are accounted for by $\frac{1-x}{1-2x}$, so assume $\pi \in B(1200)$ contains a letter greater than 1.  First suppose $\pi$ is decomposable as
\begin{equation}\label{1200form1}
\pi=0^{a_0}1^{b_1}0^{a_1}\cdots1^{b_\ell}0^{a_\ell}\pi'\sigma, \qquad \ell \geq 1,
\end{equation}
where $a_j>0$ for $0 \leq j \leq \ell-1$ with $a_\ell \geq0$, $b_j>0$ for all $j$, $\sigma=0$ or $\varepsilon$, and $\pi' \neq \varepsilon$ starts with $i \in [2,\ell+1]$ and does not contain 0. Note $\pi'$ is nondecreasing with at most $\ell+1-i$ jumps, and hence is enumerated by $h_{\ell+1-i}$. Considering all possible $\ell$ and $i$ then implies $\pi$ of the form \eqref{1200form1} are enumerated by
\begin{align*}
&\sum_{\ell \geq1}\frac{x^{2\ell}(1+x)}{(1-x)^{2\ell+1}}\sum_{i=2}^{\ell+1}\frac{x}{1-2x}\left(\frac{1-x}{1-2x}\right)^{\ell+1-i}=\frac{1+x}{1-x}\sum_{\ell \geq1}\left(\frac{x}{1-x}\right)^{2\ell}\left(\left(\frac{1-x}{1-2x}\right)^\ell-1\right)\\
&=\frac{x^2(1+x)}{1-x}\left(\frac{1}{1-3x+x^2}-\frac{1}{1-2x}\right)=\frac{x^3(1+x)}{(1-2x)(1-3x+x^2)}.
\end{align*}

Otherwise, $\pi$ is expressible as
\begin{equation}\label{1200form2}
\pi=0^{a_0}1^{b_1}0^{a_1}\cdots1^{b_\ell}0^{a_\ell}\pi'0\pi'', \qquad \ell \geq1,
\end{equation}
where the $a_i$ and $b_i$ are as before, $\pi'$ and $\pi''$ are nonempty and $\pi'$ starts with $i$ and has exactly $j$ jumps for some $i \in [2,\ell+1]$ and $0 \leq j \leq \ell+1-i$.  Note $\pi$ avoiding $1200$ implies neither $\pi'$ nor $\pi''$ can contain 0, and hence are nondecreasing.
If $p$ denotes the last letter of $\pi'$ and $q$ the first letter of $\pi''$, then we have $q=p+a$ for some $0 \leq a \leq \ell+2-i-j$, as $\pi'$ starting with $i$ and containing exactly $j$ jumps implies that at most $\ell+1-i-j$ integers can be skipped in going from $p$ to $q$.  Then $\pi''$ starting with $p=q+a$ implies it can contain at most $\ell+2-i-j-a$ jumps for each $a$.

Let $h_j^*=h_j^*(x)$ denote the generating function that counts nonempty nondecreasing sequences starting with 0 and containing exactly $j$ jumps.  Then we have $h_0^*=h_0$ and $h_j^*=h_j-h_{j-1}$ for $j \geq 1$, by the definitions.  By Lemma \ref{1000lem}, we have
$$h_j^*(x)=\frac{x}{1-2x}\left(\left(\frac{1-x}{1-2x}\right)^j-\left(\frac{1-x}{1-2x}\right)^{j-1}\right)=\frac{x^2(1-x)^{j-1}}{(1-2x)^{j+1}}, \qquad j \geq 1.$$
Considering all possible $\ell$, $i$, $j$ and $a$ implies $\pi$ of the form \eqref{1200form2} are enumerated by
$$\sum_{\ell \geq1}\frac{x^{2\ell}}{(1-x)^{2\ell+1}}\sum_{i=2}^{\ell+1}\sum_{a=0}^{\ell+2-i}xh_0(x)h_{\ell+2-i-a}(x)+\sum_{\ell \geq2}\frac{x^{2\ell}}{(1-x)^{2\ell+1}}\sum_{i=2}^{\ell}\sum_{j=1}^{\ell+1-i}h_j^*(x)\sum_{a=0}^{\ell+2-i-j}xh_{\ell+2-i-j-a}(x),$$
where we have differentiated the cases when $j=0$ or $j \geq 1$.  The first multi-sum in the preceding expression may be simplified as follows:
\begin{align*}
&\frac{x^3}{(1-x)(1-2x)^2}\sum_{\ell \geq1}\left(\frac{x}{1-x}\right)^{2\ell}\sum_{i=2}^{\ell+1}\sum_{a=0}^{\ell+2-i}\left(\frac{1-x}{1-2x}\right)^{\ell+2-i-a}\\
&=\frac{x^2}{(1-x)(1-2x)}\sum_{\ell \geq1}\left(\frac{x}{1-x}\right)^{2\ell}\sum_{i=2}^{\ell+1}\left(\left(\frac{1-x}{1-2x}\right)^{\ell+3-i}-1\right)\\
&=\frac{x(1-x)}{(1-2x)^2}\sum_{\ell \geq1}\left(\frac{x^2}{(1-x)(1-2x)}\right)^{\ell}\left(1-\left(\frac{1-2x}{1-x}\right)^\ell\right)-\frac{x^4}{(1-x)^3(1-2x)}\sum_{\ell\geq1}\ell\left(\frac{x^2}{(1-x)^2}\right)^{\ell-1}\\
&=\frac{x^4(1-x)^2}{(1-2x)^3(1-3x+x^2)}-\frac{x^4(1-x)}{(1-2x)^3}.
\end{align*}

The second multi-sum above is given by
\begin{align}
&\sum_{\ell \geq2}\frac{x^{2\ell}}{(1-x)^{2\ell+1}}\sum_{i=2}^\ell\sum_{j=1}^{\ell+1-i}\frac{x^2(1-x)^{j-1}}{(1-2x)^{j+1}}\sum_{a=0}^{\ell+2-i-j}\frac{x^2}{1-2x}\left(\frac{1-x}{1-2x}\right)^{\ell+2-i-j-a}\notag\\
&=\frac{x^4}{(1-2x)^4}\sum_{\ell \geq2}\left(\frac{x^2}{(1-x)(1-2x)}\right)^\ell\sum_{i=2}^\ell\sum_{j=1}^{\ell+1-i}\sum_{a=0}^{\ell+2-i-j}\left(\frac{1-2x}{1-x}\right)^{i+a}.\label{multsum2*}
\end{align}
Simplifying the inner triple sum in the last expression, we have
\begin{align*}
&\sum_{i=2}^\ell\sum_{j=1}^{\ell+1-i}\sum_{a=0}^{\ell+2-i-j}\left(\frac{1-2x}{1-x}\right)^{i+a}=\frac{1-x}{x}\sum_{i=2}^\ell\left(\frac{1-2x}{1-x}\right)^i\sum_{j=1}^{\ell+1-i}\left(1-\left(\frac{1-2x}{1-x}\right)^{\ell+3-i-j}\right)\\
&=\frac{1-x}{x}\sum_{i=2}^{\ell}(\ell+1-i)\left(\frac{1-2x}{1-x}\right)^i-\frac{1-x}{x}\sum_{i=2}^\ell\left(\frac{1-2x}{1-x}\right)^{\ell+3}\sum_{j=1}^{\ell+1-i}\left(\frac{1-x}{1-2x}\right)^j\\
&=\frac{1-x}{x}\sum_{i=1}^{\ell-1}i\left(\frac{1-2x}{1-x}\right)^{\ell+1-i}+\frac{(1-x)^2}{x^2}\sum_{i=2}^{\ell+1}\left(\left(\frac{1-2x}{1-x}\right)^{\ell+3}-\left(\frac{1-2x}{1-x}\right)^{i+2}\right).
\end{align*}
Note that the last formula yields zero when $\ell=1$, and hence one may start the outer sum in $\ell$ on the right side of \eqref{multsum2*} at $\ell=1$.  Substituting this last expression for the inner triple sum into \eqref{multsum2*}, we get
\begin{align*}
&\frac{x^3(1-x)}{(1-2x)^4}\sum_{\ell \geq2}\left(\frac{x}{1-x}\right)^{2\ell}\sum_{i=1}^{\ell-1}i\left(\frac{1-x}{1-2x}\right)^{i-1}+\frac{x^2}{(1-x)(1-2x)}\sum_{\ell\geq1}\ell\left(\frac{x}{1-x}\right)^{2\ell}\\
&\quad-\frac{x^2}{(1-x)^2}\sum_{\ell\geq1}\left(\frac{x^2}{(1-x)(1-2x)}\right)^\ell\sum_{i=2}^{\ell+1}\left(\frac{1-2x}{1-x}\right)^{i-2}\\
&=\frac{x^3(1-x)}{(1-2x)^4}\sum_{i\geq1}i\left(\frac{1-x}{1-2x}\right)^{i-1}\cdot\frac{(x^2/(1-x)^2)^{i+1}}{1-\frac{x^2}{(1-x)^2}}+\frac{x^2}{(1-x)(1-2x)}\cdot\frac{x^2}{(1-x)^2\left(1-\frac{x^2}{(1-x)^2}\right)^2}\\
&\quad-\frac{x}{1-x}\sum_{\ell \geq1}\left(\frac{x^2}{(1-x)(1-2x)}\right)^\ell\left(1-\left(\frac{1-2x}{1-x}\right)^\ell\right)\\
&=\frac{x^7}{(1-x)(1-2x)^5}\sum_{i\geq1}i\left(\frac{x^2}{(1-x)(1-2x)}\right)^{i-1}+\frac{x^4(1-x)}{(1-2x)^3}-\frac{x^3}{1-x}\left(\frac{1}{1-3x+x^2}-\frac{1}{1-2x}\right)\\
&=\frac{x^7(1-x)}{(1-2x)^3(1-3x+x^2)^2}+\frac{x^4(1-x)}{(1-2x)^3}-\frac{x^4}{(1-2x)(1-3x+x^2)}.
\end{align*}

Thus, from the expressions found above for the two multi-sums, we have that $\pi$ of the form \eqref{1200form2} are enumerated by
$$\frac{x^4(1-x)^2}{(1-2x)^3(1-3x+x^2)}+\frac{x^7(1-x)}{(1-2x)^3(1-3x+x^2)^2}-\frac{x^4}{(1-2x)(1-3x+x^2)}.$$
Combining the prior cases of $\pi$, and observing some cancellation, we get
\begin{align*}
f_{1200}&=\frac{1-x}{1-2x}+\frac{x^3}{(1-2x)(1-3x+x^2)}+\frac{x^4(1-x)^2}{(1-2x)^3(1-3x+x^2)}+\frac{x^7(1-x)}{(1-2x)^3(1-3x+x^2)^2}\\
&=\frac{1-x}{1-2x}+\frac{x^3}{(1-3x+x^2)^2},
\end{align*}
as before.

To complete the proof of \eqref{1200the1}, we define a bijection $g$ between $B_n(1200)$ and $B_n(1220)$.  Let $\pi \in B_n(1200)$.  If $\pi$ is binary or does not contain a 0 to the right of its first letter greater than 1, then let $g(\pi)=\pi$.  Otherwise, consider the nondecreasing section $\pi'$ within $\pi$ of the form \eqref{1200form1} or \eqref{1200form2}.  We replace each (maximal) string $w^b$ in $\pi'$, where $w \geq 2$ and $b \geq 1$, with $w0^{b-1}$, and leave all other letters of $\pi$ undisturbed.  Let $g(\pi)$ denote the resulting ascent sequence, which is seen to avoid $1220$.  Combining the two cases above, one may verify that $g$ provides the desired bijection between $B_n(1200)$ and $B_n(1220)$, which completes the proof.
\end{proof}

\end{document}
